\documentclass[11pt, letterpaper]{article}

% -----------------------------------------------------------------------
% Packages
% -----------------------------------------------------------------------
\usepackage[utf8]{inputenc}
\usepackage[T1]{fontenc}
\usepackage[expansion=false]{microtype}       % better text rendering
\usepackage{amsmath, amsfonts, amssymb}
\usepackage{geometry}
\usepackage{booktabs}
\usepackage{array}
\usepackage{xcolor}
\usepackage{hyperref}
\usepackage{enumitem}
\usepackage{titlesec}
\usepackage{parskip}
\usepackage{fancyhdr}
\usepackage{abstract}
\usepackage{listings}
\usepackage[linesnumbered, ruled, vlined]{algorithm2e}
\usepackage{caption}
\usepackage{subcaption}

% -----------------------------------------------------------------------
% Layout
% -----------------------------------------------------------------------
\geometry{margin=1in, top=1.1in, bottom=1.1in}

\hypersetup{
  colorlinks = true,
  linkcolor  = blue!70!black,
  urlcolor   = cyan!70!black,
  citecolor  = green!50!black
}

% Section formatting — small caps title, rule underneath
\titleformat{\section}
  {\large\bfseries\scshape}
  {\thesection.}{0.6em}{}[\vspace{-0.4em}\rule{\linewidth}{0.4pt}]

\titleformat{\subsection}
  {\normalsize\bfseries}
  {\thesubsection.}{0.5em}{}

\titleformat{\subsubsection}
  {\normalsize\itshape}
  {\thesubsubsection.}{0.5em}{}

% -----------------------------------------------------------------------
% Code listing style
% -----------------------------------------------------------------------
\lstdefinestyle{cstyle}{
  language        = C,
  basicstyle      = \ttfamily\small,
  keywordstyle    = \color{blue!70!black}\bfseries,
  commentstyle    = \color{gray}\itshape,
  stringstyle     = \color{orange!80!black},
  numbers         = left,
  numberstyle     = \tiny\color{gray},
  numbersep       = 8pt,
  frame           = single,
  frameround      = tttt,
  rulecolor       = \color{black!20},
  breaklines      = true,
  captionpos      = b,
  tabsize         = 4,
  showstringspaces= false
}
\lstset{style=cstyle}

% -----------------------------------------------------------------------
% Header / Footer
% -----------------------------------------------------------------------
\pagestyle{fancy}
\fancyhf{}
\rhead{\small\scshape GA/TSP Technical Specification}
\lhead{\small Joel Maldonado --- University of Arizona}
\rfoot{\small\thepage}
\renewcommand{\headrulewidth}{0.4pt}

% -----------------------------------------------------------------------
% Abstract formatting
% -----------------------------------------------------------------------
\renewcommand{\abstractnamefont}{\bfseries\scshape}
\renewcommand{\abstracttextfont}{\small}
\setlength{\absleftindent}{0.5in}
\setlength{\absrightindent}{0.5in}

% -----------------------------------------------------------------------
% Convenience macros
% -----------------------------------------------------------------------
\newcommand{\hpcnote}[1]{%
  \vspace{0.4em}
  \noindent\colorbox{blue!6}{\parbox{\dimexpr\linewidth-2\fboxsep}{%
    \footnotesize\textbf{HPC Note:}\enspace #1
  }}
  \vspace{0.4em}
}

\newcommand{\warn}[1]{%
  \vspace{0.4em}
  \noindent\colorbox{orange!10}{\parbox{\dimexpr\linewidth-2\fboxsep}{%
    \footnotesize\textbf{GPU Constraint:}\enspace #1
  }}
  \vspace{0.4em}
}

% -----------------------------------------------------------------------
% Document
% -----------------------------------------------------------------------
\title{%
  \textbf{Technical Specification: GPGPU-Ready Genetic Algorithm}\\[0.4em]
  \large Architecture for the Symmetric Traveling Salesman Problem
}
\author{%
  Joel Maldonado\\
  \small Department of Applied Mathematics, University of Arizona\\
  \small \href{mailto:joelmaldonado@arizona.edu}{joelmaldonado@arizona.edu}
}
\date{April 2026 \quad|\quad Version 1.0}

\begin{document}

\maketitle
\thispagestyle{fancy}

\begin{abstract}
This document specifies the software architecture for a modular, sequential Genetic
Algorithm (GA) targeting the Traveling Salesman Problem (TSP), implemented in
standard C99.  Although the initial execution target is a single CPU core, every
architectural decision is made with NVIDIA CUDA in mind: memory layouts prefer
coalescing patterns, all functions are stateless with respect to global storage,
random-number generators are per-individual, and evaluation is batch-oriented to
map directly onto kernel launches.  The document covers the formal problem
definition, layer decomposition, data-structure design with AoS/SoA analysis,
genetic operator specifications, a formal algorithm statement, complexity
characterization, and a checklist of CUDA-readiness constraints.
\end{abstract}

\tableofcontents
\bigskip

% ======================================================================
\section{Problem Formalization}
% ======================================================================

The Symmetric Traveling Salesman Problem is defined on a complete weighted graph
$G = (V, E)$ where $|V| = n$.  The distance matrix $D \in \mathbb{R}^{n \times n}$
satisfies $d_{ij} = d_{ji} \geq 0$ and $d_{ii} = 0$.  Given the coordinates
$(x_k, y_k)$ for city $k$, the Euclidean distance is
\[
  d_{ij} = \sqrt{(x_i - x_j)^2 + (y_i - y_j)^2}.
\]
A \emph{tour} is a bijection $\pi: \{0,\dots,n-1\} \to \{0,\dots,n-1\}$ encoding
a Hamiltonian cycle.  The objective is
\begin{equation}
  \min_{\pi \in S_n}\; L(\pi)
  = \sum_{k=0}^{n-2} d_{\pi_k,\,\pi_{k+1}} + d_{\pi_{n-1},\,\pi_0},
  \label{eq:objective}
\end{equation}
where $S_n$ denotes the symmetric group on $n$ elements.  Because $L$ is
symmetric in $d$, the search space has effective size $(n-1)!/2$.

% ======================================================================
\section{System Architecture}
% ======================================================================

The implementation is decomposed into five loosely-coupled layers.  Each layer
exposes a well-defined functional interface and holds no mutable global state,
a constraint that directly enables CUDA kernel extraction.

\begin{table}[h]
\centering
\caption{Layer definitions and forward-looking GPU porting strategies.}
\label{tab:layers}
\renewcommand{\arraystretch}{1.35}
\begin{tabular}{@{} p{2.4cm} p{4.2cm} p{5.8cm} @{}}
\toprule
\textbf{Layer} & \textbf{Responsibility} & \textbf{CUDA Porting Strategy} \\
\midrule
Instance        & Coordinate I/O; distance matrix construction
                & Pin $D$ to constant memory ($n \leq 256$) or L2-cached global memory \\
Representation  & Population and tour memory management
                & Migrate AoS layout to SoA for coalesced access \\
Evaluation      & Fitness computation; tour validation
                & One thread per individual; parallel reduction within tour \\
Evolution       & Selection, crossover, mutation, elitism
                & Warp-level primitives (\texttt{\_\_shfl\_sync}) for tournament logic \\
Control         & GA main loop; termination criteria
                & Remains host-side; issues asynchronous kernel launches \\
\bottomrule
\end{tabular}
\end{table}

% ======================================================================
\section{Data Structures and Memory Layout}
% ======================================================================

\subsection{Distance Matrix}

The matrix $D$ is allocated once at startup and treated as read-only for the
entirety of GA execution.  It is stored in row-major order as a flat array of
\texttt{double}, indexed by $i \cdot n + j$.

\hpcnote{For small instances ($n \leq 256$, i.e., $\leq 512$\,KiB at 64-bit),
$D$ fits in CUDA constant memory and benefits from broadcast reads when all
threads in a warp address the same entry.  For larger instances it must reside
in global memory; a symmetric half-matrix layout ($n(n+1)/2$ elements) reduces
footprint and may improve cache utilization.}

\subsection{Tour Representation}

A tour is a permutation stored as an integer array of length $n$.  The
associated struct carries precomputed fitness to avoid redundant evaluation:

\begin{lstlisting}[caption={Sequential tour struct (AoS element).}]
typedef struct {
    int    *cities;    /* permutation of {0, ..., n-1}         */
    double  length;    /* total Euclidean tour length           */
    double  fitness;   /* derived from length; higher is better */
} Tour;
\end{lstlisting}

\textbf{Invariant.}  After any genetic operator, every valid tour $\pi$ must
satisfy $\{\pi_0, \dots, \pi_{n-1}\} = \{0, \dots, n-1\}$.  Violations
indicate a logic error and are caught by the validation layer
(Section~\ref{sec:validation}).

\subsection{Population Layout: AoS vs.\ SoA}
\label{sec:layout}

The sequential implementation uses an \emph{Array of Structures} (AoS), which
is natural for single-threaded access patterns.

\begin{lstlisting}[caption={AoS: simple but GPU-hostile.}]
Tour population[POP_SIZE];
/* Thread i accesses population[i].cities[k]:
   addresses stride by sizeof(Tour) -- non-coalesced. */
\end{lstlisting}

\warn{On a GPU, when thread $i$ and thread $i+1$ access \texttt{population[i].cities[k]}
and \texttt{population[i+1].cities[k]} respectively, the hardware must issue
$W$ independent 32-byte transactions for a warp of $W=32$ threads because the
\texttt{cities} pointers are separated by \texttt{sizeof(Tour)} bytes.  This
fully serializes what should be a single coalesced load.}

The GPU-ready layout separates gene arrays from metadata:

\begin{lstlisting}[caption={SoA: GPU-friendly, migration target.}]
/* All gene-k values for the entire population are contiguous. */
int    *pop_cities;    /* [POP_SIZE * n]  row i: cities of individual i  */
double *pop_length;    /* [POP_SIZE]                                      */
double *pop_fitness;   /* [POP_SIZE]                                      */

/* Thread i reads gene k: pop_cities[i*n + k].
   Threads 0..31 reading k=0 hit addresses {0,n,2n,...,31n} --
   still strided.  Transpose to [n * POP_SIZE] with layout
   pop_cities[k*POP_SIZE + i] for fully coalesced access. */
\end{lstlisting}

The transposed SoA layout (\texttt{pop\_cities[k * POP\_SIZE + i]}) ensures
that a warp reading gene $k$ of all individuals issues a single coalesced
128-byte transaction, reducing memory traffic by up to $32\times$.

% ======================================================================
\section{Genetic Operators}
% ======================================================================

\subsection{Initialization}

The initial population is generated by Fisher--Yates shuffles seeded from a
per-individual RNG state.  Diversity is guaranteed by construction; no
duplicate-rejection loop is needed.

\subsection{Fitness Evaluation}

Tour length is computed directly from Equation~\eqref{eq:objective} using the
precomputed matrix $D$.  Fitness is defined as the inverse of length,
$f(\pi) = 1/L(\pi)$, so that higher fitness corresponds to shorter tours.

\textbf{Parallelization structure.}  Evaluation is \emph{embarrassingly
parallel} across the population: individual $i$'s fitness depends only on its
own permutation and the read-only matrix $D$.  A future CUDA kernel assigns
one thread block per individual, with each thread summing a contiguous segment
of edges and a final warp reduction collecting the total.

\subsection{Selection: Binary Tournament}

For each offspring slot, two individuals are drawn uniformly at random and the
fitter one is selected.  Tournament size $k=2$ is the default; $k$ can be
tuned to control selection pressure without restructuring the algorithm.

\hpcnote{Tournament selection requires no global information: each contest is
$O(k)$ work and completely independent.  This contrasts with fitness-proportionate
(roulette-wheel) selection, which requires an $O(N \log N)$ prefix scan over
the entire population --- a bottleneck even with GPU-accelerated prefix sums.}

\subsection{Crossover: Order Crossover (OX1)}
\label{sec:ox1}

Order Crossover preserves the relative ordering of cities inherited from one
parent while filling remaining positions from the other.  The algorithm proceeds
in three steps: (1) copy a contiguous segment $[\ell, r]$ from parent~$A$
directly into the offspring; (2) traverse parent~$B$ starting at position
$r+1 \pmod{n}$ and append cities not already present; (3) wrap to fill all
positions.

The sequential implementation uses an $O(n)$ boolean lookup buffer to test
membership.  The planned GPU implementation replaces this with a 64-bit
bitmask held in registers (valid for $n \leq 64$) or shared memory (valid for
$n \leq 1024$), eliminating branch divergence caused by cache-miss-dependent
conditional branches in the linear scan.

\warn{OX1 contains data-dependent branching proportional to $n$, which can
cause warp divergence when different threads reach different branch outcomes at
the same instruction.  Mitigation: predicated instructions and bitmask
membership tests that resolve without branching.}

\subsection{Mutation: Swap and Inversion}

Two mutation operators are supported:

\begin{itemize}[noitemsep]
  \item \textbf{Swap}: select two positions $i \neq j$ uniformly at random;
    exchange $\pi_i \leftrightarrow \pi_j$.  Complexity: $O(1)$.
  \item \textbf{Inversion (2-opt move)}: select a sub-segment $[\ell, r]$;
    reverse it in place.  Complexity: $O(n)$.
\end{itemize}

Each operator is applied independently per individual; there is no
inter-individual communication, making mutation trivially parallel.

\warn{Mixing operator types within a single kernel launch introduces warp
divergence if different threads branch to different operators.  Prefer separate
kernel launches per operator type, or compile a single-operator kernel and
execute them sequentially on the host.}

\subsection{Elitism}

The top $e$ individuals (by fitness) are copied unchanged into the next
generation before new offspring are generated.  This prevents fitness
regression across generations.

For the sequential baseline, elitism is implemented via a partial sort
(introselect) in $O(N)$ expected time.  On the GPU, a parallel Top-$e$
selection via bitonic sort on the fitness array is preferred over a full sort
to minimize kernel occupancy.

% ======================================================================
\section{Algorithm Statement}
% ======================================================================

\begin{algorithm}[H]
\DontPrintSemicolon
\SetAlgoLined
\LinesNumbered
\KwIn{TSP instance $I$ (coordinates, $n$); GA configuration $C$
      (population size $N$, elite count $e$, tournament size $k$,
      crossover rate $p_c$, mutation rate $p_m$, maximum generations $G$)}
\KwOut{Best tour $\pi^*$ found}
\BlankLine
$D \gets \textsc{BuildDistanceMatrix}(I)$\;
$P \gets \textsc{InitializePopulation}(N, n)$\;
$\textsc{EvaluateFitness}(P,\, D)$\;
\BlankLine
\For{$g \gets 1$ \KwTo $G$}{
  $P' \gets \textsc{ExtractElites}(P,\, e)$\;
  \While{$|P'| < N$}{
    $a \gets \textsc{Tournament}(P,\, k)$\;
    $b \gets \textsc{Tournament}(P,\, k)$\;
    \eIf{$\textsc{Uniform}(0,1) < p_c$}{
      $c \gets \textsc{OX1Crossover}(a,\, b)$\;
    }{
      $c \gets \textsc{Copy}(a)$\;
    }
    \If{$\textsc{Uniform}(0,1) < p_m$}{
      $\textsc{Mutate}(c)$\;
    }
    $\textsc{Append}(P',\, c)$\;
  }
  $P \gets P'$\;
  $\textsc{EvaluateFitness}(P,\, D)$\;
  $\textsc{LogStatistics}(g,\, P)$\;
}
\BlankLine
\Return $\arg\max_{\pi \in P} f(\pi)$\;
\caption{Sequential Reference GA for TSP}
\label{alg:ga}
\end{algorithm}

% ======================================================================
\section{Complexity Analysis}
% ======================================================================

Let $N$ denote population size, $n$ the number of cities, $G$ the number of
generations, $e$ the elite count, and $k$ the tournament size.

\begin{table}[h]
\centering
\caption{Per-generation computational complexity.}
\label{tab:complexity}
\renewcommand{\arraystretch}{1.35}
\begin{tabular}{@{} l l l @{}}
\toprule
\textbf{Operation} & \textbf{Sequential} & \textbf{Parallelizable?} \\
\midrule
Distance matrix build  & $O(n^2)$ once   & Yes, $O(n)$ with $n$ threads \\
Fitness evaluation     & $O(Nn)$         & Yes, $O(n)$ with $N$ threads \\
Tournament selection   & $O(kN)$         & Yes, $O(k)$ with $N$ threads \\
OX1 crossover          & $O(Nn)$         & Yes, $O(n)$ with $N$ threads \\
Swap mutation          & $O(N)$          & Yes, $O(1)$ with $N$ threads \\
Inversion mutation     & $O(Nn)$         & Yes, $O(n)$ with $N$ threads \\
Elite extraction       & $O(N)$          & Partial (bitonic sort) \\
\midrule
\textbf{Total per generation} & $O(GNn)$ & Reduces to $O(Gn)$ on GPU \\
\bottomrule
\end{tabular}
\end{table}

The total wall-clock complexity for the sequential implementation is $O(G N n)$.
A GPU implementation with $N$ threads in parallel reduces the dominant
evaluation and crossover terms from $O(Nn)$ to $O(n)$ per generation,
yielding a theoretical $N$-fold speedup on those kernels.

% ======================================================================
\section{CUDA Readiness Constraints}
% ======================================================================

The following constraints are enforced in the sequential baseline to ensure
that GPU porting is an incremental refactor rather than a rewrite.

\begin{enumerate}[label=\textbf{CR-\arabic*.}, leftmargin=*, itemsep=0.5em]

  \item \textbf{Zero global mutable state.}
    No \texttt{static}, \texttt{extern}, or file-scope mutable variables are
    used.  All data flows through explicit pointer arguments.  This is a
    hard requirement because CUDA device functions have no access to host
    address space.

  \item \textbf{Per-individual RNG state.}
    Each individual carries its own PRNG state (a \texttt{uint64\_t} for a
    splitmix64 generator).  On the GPU, this maps directly to a
    \texttt{curandState} per thread, avoiding the serialization caused by a
    single shared generator.

  \item \textbf{Batch-oriented evaluation interface.}
    The evaluation function signature is
    \texttt{evaluate\_population(Tour*, int N, double* D, int n)}, not a loop
    over \texttt{evaluate\_tour} calls.  The batch signature mirrors a CUDA
    kernel launch: the host passes a device pointer and the GPU fills results
    in parallel.

  \item \textbf{Flat, pointer-based memory.}
    No nested dynamic allocation (\texttt{Tour**} with per-element
    \texttt{malloc}).  All population data is held in contiguous flat buffers
    compatible with \texttt{cudaMalloc} and \texttt{cudaMemcpy}.

  \item \textbf{Deterministic seeding.}
    A command-line seed fully determines all random choices, enabling exact
    reproduction of any run.  This is critical for isolating GPU race
    conditions: compare GPU output against the seeded sequential baseline.

  \item \textbf{SoA migration path documented.}
    The AoS-to-SoA transformation described in Section~\ref{sec:layout} is
    the first step of the CUDA port.  No other algorithmic change is required
    before writing kernel code.

\end{enumerate}

% ======================================================================
\section{Validation Layer}
\label{sec:validation}
% ======================================================================

Correctness of permutation operators is verified by a dedicated validator
executed after every crossover and mutation in debug builds (\texttt{\#ifdef
DEBUG}).

\textbf{Permutation check.}  An $O(n)$ pass marks each city encountered in
the tour; any duplicate or missing index raises an assertion:

\begin{lstlisting}[caption={Tour validation.}]
bool validate_tour(const Tour *t, int n) {
    bool seen[n];
    memset(seen, 0, n * sizeof(bool));
    for (int k = 0; k < n; k++) {
        if (t->cities[k] < 0 || t->cities[k] >= n) return false;
        if (seen[t->cities[k]])                      return false;
        seen[t->cities[k]] = true;
    }
    return true;
}
\end{lstlisting}

\textbf{Fitness consistency check.}  After evaluation, $L(\pi)$ is
independently recomputed and compared against the stored value within a
relative tolerance of $10^{-9}$.

\textbf{Relevance to GPU debugging.}  The validator provides a ground-truth
reference against which GPU kernel output can be checked element-wise.
Memory corruption from out-of-bounds writes in a CUDA kernel typically
manifests as a permutation invariant violation before causing a numerical
anomaly, making the validator an early-warning system for kernel bugs.

% ======================================================================
\section{Summary}
% ======================================================================

This specification defines a sequential GA for the TSP whose architecture is
shaped by the eventual requirements of a CUDA implementation.  The key design
decisions are:

\begin{itemize}[noitemsep]
  \item \textbf{Stateless functions} to enable direct kernel extraction.
  \item \textbf{Flat contiguous memory} compatible with device allocation.
  \item \textbf{Batch evaluation interface} mirroring kernel launch semantics.
  \item \textbf{Per-individual RNG} ready for \texttt{cuRAND} substitution.
  \item \textbf{Documented AoS-to-SoA transition} as the primary migration
    step.
  \item \textbf{Tournament selection} to avoid global prefix-sum dependencies.
\end{itemize}

The result is a CPU baseline that can be validated rigorously and then ported
to GPU by replacing the evaluation and evolution loops with kernel launches ---
no structural redesign required.

\end{document}